\section{Contribution, limitations and further questions}

The mixed-product theorem identifies an attained maximum when there is no inconsistent centered singleton. Its proof keeps the exact depleted moment, merges colors through positive partition weights, and handles arbitrary multiplicities through affine indicator identities. The scalar consequences recover the spectral estimate with a smaller singleton factor and an optimal universal prefactor. The diagonal results determine the original cubic threshold, the parity transition, a hierarchy of further corrections and the sharp first uniform error.

The base recurrence, its normalized parity form and the leading coefficients belong to Collins and Nagatsu \cite{CN}. The mixed estimate itself is already part of Bordenave's spectral-gap method \cite{Bordenave}. The special-function identities and the saddle method are classical. The precise refinements proved here are the attained envelope at fixed support and singleton counts, the necessity and bounded-comparability characterization of the cubic scale, the uniform sublinear correction hierarchy, and the sharp parity-dependent error on bounded square-root ranges. A targeted literature check found no exact duplicate of these statements, but cannot certify the absence of every earlier equivalent formula.

Bordenave and Zhu \cite[Lemma 2.1]{BZ} retain the original singleton factor $3k/\sqrt N$ and explicitly bound the raw factorial ratio by $\sqrt e$ in their equation (2.3). That elementary bound is a known ingredient here. Our exact refinement retains $\mu_h(n,a-h)$ in addition to $F_n(a-h)$. The factor $\Gamma_{n,a,h}$ is a scalar upper bound on that moment and is smaller than $\beta$ when $h<a$. The attained-maximum assertion concerns $b=0$. We do not claim that $2^b$ is sharp for prescribed $b>0$.

The Laguerre representation makes the parameter regime precise: the degree is $k$, the parameter is $\alpha=-N-1$, and the argument is $-N$. Thus $-\alpha/k\to\infty$ throughout the sublinear range. Kuijlaars and McLaughlin \cite{KM}, and the regime described by Dai and Wong \cite{DW}, treat a finite limiting ratio greater than one. Those results overlap the macroscopic part of our full-range formula, so that part is not presented as a new special-function asymptotic regime. Theorem 1 of Dai, Ismail and Wang \cite{DIW} assumes $\alpha>0$ when $\alpha/k\to\infty$. In the inspected preprint of Dunster, Gil and Segura \cite[Section 1, equation (1.6) and Cases 1a, 1b and 2]{DGS}, $\alpha=(k+1/2)(a^2-1)$ with $a\ge0$, excluding $\alpha=-N-1$ for $k\le N$. These parameter comparisons do not turn analytic continuation into a uniform error estimate.

At $k=N$, the identity with a truncated exponential is classical, as recorded in \cite[Section 1]{DMO}. The contribution of Theorem~\ref{thm:global} is a self-contained formula that patches fixed orders, the parity transition and the endpoint while keeping $(N)_k$ exact. Its stated $O(N^{-1/4})$ global bound is not asserted to be optimal. The sharp rate in Section~\ref{sec:rates} concerns the different, quartically corrected hyperbolic approximation.

The inequalities are proved symbolically for their stated parameter ranges. Exact finite computations check normalization, the depleted recurrence and the conditioning examples. Numerical plots illustrate the asymptotic regimes. Neither kind of computation is used as a substitute for a general proof.

The remaining graph-level task is to combine the sharper cell-pattern estimate with circuit enumeration and tangle remainders while retaining the degree dependence. A better estimate for one fixed product does not automatically improve a spectral theorem: the patterns that dominate the full sum must be identified first. The present paper does not assert a new second-eigenvalue bound or an end-to-end audit of either random-graph proof.
